\chapter{Table des notations}

% MIGRÉ depuis memoire/main.tex (plages 3230-3259).
% RELECTURE FAITE (août 2026) : la table héritée ne portait AUCUN symbole du
% modèle de cascade (W, S, A, TRANS, ROOT, g, s_j, e_j, R_0, t, kappa), alors
% que c'est le modèle central. Elle est réécrite et organisée en cinq blocs.
\label{not:ann:notations}

\noindent Les symboles sont regroupés par brique du modèle, dans l'ordre où le mémoire les
introduit. La colonne de droite renvoie au chapitre qui les définit.

\renewcommand{\arraystretch}{1.22}
\begin{center}\small
\begin{tabular}{>{\raggedright}p{3.3cm} >{\raggedright\arraybackslash}p{10.2cm} c}
\toprule
\rowcolor{lightblue}
\textbf{Symbole} & \textbf{Signification} & \textbf{Ch.} \\
\midrule
\multicolumn{3}{@{}l}{\textit{\textcolor{navy}{Piliers et cascade dirigée}}} \\
$P_1,\dots,P_5$ & Les cinq piliers DORA : gouvernance, gestion des incidents, tests de résilience, prestataires tiers, partage d'informations & \ref{chap:cascade} \\
$\mathrm{TRANS}$ & Matrice d'influence brute du classeur qualitatif, asymétrique, non normalisée & \ref{chap:cascade} \\
$\mathrm{ROOT}$ & Propension de chaque pilier à \emph{amorcer} une chaîne & \ref{chap:cascade} \\
$s_j$ & Somme de la ligne $j$ de $\mathrm{TRANS}$ ; $s_{\max}=\max_j s_j$ & \ref{chap:cascade} \\
$g$ & Gain de propagation, part de contagion du pilier le plus exposé & \ref{chap:cascade} \\
$W$ & Matrice de contagion normalisée à la Leontief, $W=g\,\mathrm{TRANS}/s_{\max}$ & \ref{chap:cascade} \\
$e_j$ & Propension à propager depuis $j$, $e_j=g\,s_j/s_{\max}$ & \ref{chap:cascade} \\
$\rho(W)$ & Rayon spectral de $W$ ; $\rho(W)\le g<1$ garantit la stabilité & \ref{chap:cascade} \\
$(I-W)^{-1}$ & Inverse de Leontief, descendance espérée avec multiplicité & \ref{chap:cascade} \\
$M,\ R_0$ & Matrice de génération suivante et son rayon spectral (extinction si $R_0<1$) & \ref{chap:cascade} \\
$S_a$ & Ensemble des piliers atteints depuis l'amorce $a$ & \ref{chap:cascade} \\
$K_j,\ u_j$ & Seuil sur la latente ; seuil de matérialité observable (art.~18) & \ref{chap:cascade} \\
\midrule
\multicolumn{3}{@{}l}{\textit{\textcolor{navy}{Identifiabilité et identification partielle}}} \\
$S,\ A$ & Parties symétrique et antisymétrique de $W$ : co-occurrence (identifiée) et direction (non identifiée) & \ref{ch:identifiabilite} \\
$J_{jk}$ & Courant de probabilité $\pi_jP_{jk}-\pi_kP_{kj}$ ; nul ssi réversibilité & \ref{ch:identifiabilite} \\
$\pi,\ \tilde P$ & Loi stationnaire ; noyau adjoint (chaîne renversée dans le temps) & \ref{ch:identifiabilite} \\
$\mathcal{E}(f,f)$ & Forme de Dirichlet, variation quadratique de la martingale ; ne dépend que de $S$ & \ref{ch:identifiabilite} \\
$\sigma$ & Production d'entropie ; nulle ssi la chaîne est réversible & \ref{ch:identifiabilite} \\
$T(M),\ z$ & Statistique d'asymétrie et son écart réduit au placebo par permutation & \ref{ch:identifiabilite} \\
$t$ & Degré d'ignorance directionnelle, $|A_{jk}|\le t\,S_{jk}$ ; $t=1$ = ignorance totale & \ref{chap:partielle} \\
$\mathcal{W}(t)$ & Ensemble admissible des matrices compatibles avec la donnée & \ref{chap:partielle} \\
$u_{ij}$ & Excès de propagation propre à la source, $\operatorname{logit}p_{ij}=\operatorname{logit}p_j+u_{ij}$ & \ref{chap:partielle} \\
$L(\mathcal{V}),\ V(\mathcal{I})$ & Largeur de bande de capital ; valeur d'une information & \ref{ch:livrable} \\
\midrule
\multicolumn{3}{@{}l}{\textit{\textcolor{navy}{Socle : fréquence et sévérité}}} \\
$\xi,\ \sigma_{\text{GPD}},\ u$ & Indice de queue, échelle et seuil POT de la GPD de sévérité & \ref{chap:socle} \\
$\zeta_u=p_u$ & Probabilité de dépassement du seuil, $\Prob(X>u)$ & \ref{chap:socle} \\
$N_u$ & Nombre d'excès au-delà de $u$ & \ref{chap:socle} \\
$\lambda_{ref}$ & Fréquence de référence PRC ($341$/an) : \emph{clé de répartition} par vecteur, n'alimente pas le moteur & \ref{chap:donnees} \\
$\lambda$ & Fréquence d'entrée du moteur, incidents \textsc{tic} matériels ($21{,}6$/an au secteur) & \ref{chap:donnees} \\
$\varphi,\ a$ & Facteur de surdispersion ; charge systémique commune du mélange lognormal & \ref{chap:socle} \\
$r,\ p$ & Paramètres de la loi binomiale négative de fréquence & \ref{chap:socle} \\
$b$ & Élasticité de la sévérité à la taille de la firme & \ref{chap:donnees} \\
\midrule
\multicolumn{3}{@{}l}{\textit{\textcolor{navy}{Conformité multi-états}}} \\
$C^\star_j,\ \Theta$ & Capacité de conformité latente du pilier $j$ ; facteur systémique commun & \ref{chap:multietats} \\
$\gamma$ & Charge systémique de la latente de conformité ($0{,}68$) & \ref{chap:multietats} \\
$K_{\text{bas}},K_{\text{haut}}$ & Seuils ordonnés découpant les états NC / PC / C & \ref{chap:multietats} \\
$Q$ & Générateur de la chaîne de Markov de mise en conformité (séjours de type phase) & \ref{chap:multietats} \\
$S_0,S_1,S_2$ & Scénarios conforme / partiel / non conforme (multiplicateurs de fréquence) & \ref{chap:socle} \\
\midrule
\multicolumn{3}{@{}l}{\textit{\textcolor{navy}{Capital, allocation et décision}}} \\
$\VaR_\alpha,\ \TVaR_\alpha$ & Value-at-Risk et Tail-VaR au niveau $\alpha=99{,}5\,\%$ & \ref{chap:socle} \\
$\Delta_{DORA}$ & Surcoût de capital imputable à la non-conformité & \ref{chap:resultats} \\
$AC_k,\ \phi_j$ & Contribution d'Euler ; valeur de Shapley du pilier $j$ & \ref{chap:resultats} \\
$\mathrm{CoC}$ & Taux de coût du capital de la marge de risque ($6\,\%$, révisé à $4{,}75\,\%$) & \ref{chap:resultats} \\
$L,\ L^\star$ & Portée d'un traité en excès de perte annuel ; portée optimale & \ref{chap:resultats} \\
$\kappa$ & Coefficient d'inefficacité du transfert (capacité, base, contrepartie) & \ref{chap:resultats} \\
$\theta$ & Paramètre de la copule de Gumbel, \emph{utilisée comme axe de comparaison} et non comme modèle & \ref{ch:robustesse} \\
\bottomrule
\end{tabular}
\end{center}

\clearpage

% =====================================================================
